% !TeX program = pdflatex
% =============================================================================
% Lektionsplan (Diskussionsgrundlage, NICHT das Arbeitsblatt selbst)
% Kraftaddition & Kraftzerlegung -- UZH Praktikum I, Sara Romer Urban, KS im Lee, HS2026
% Klasse: 2e
% =============================================================================
\documentclass[11pt,a4paper]{article}

\usepackage[T1]{fontenc}
\usepackage[utf8]{inputenc}
\usepackage[ngerman]{babel}
\usepackage{lmodern}
\usepackage[a4paper,margin=2.2cm]{geometry}
\usepackage{array,booktabs,tabularx,xltabular}
\usepackage{enumitem}
\usepackage{xcolor}
\usepackage{amsmath}
\usepackage{siunitx}
\sisetup{per-mode=fraction,fraction-command=\frac}
\usepackage{url}
\Urlmuskip=0mu plus 1mu
\def\UrlBreaks{\do\/\do-\do_}

\definecolor{titleblue}{HTML}{183A6B}
\definecolor{accentblue}{HTML}{2E6F9E}
\definecolor{groupteal}{HTML}{1B7F72}
\definecolor{darkgray}{HTML}{4A4A4A}
\definecolor{warnred}{HTML}{9E2E2E}

\setlength{\parindent}{0pt}
\setlength{\parskip}{0.5em}

\newcommand{\Methode}[1]{{\small\color{groupteal}\textbf{#1}}}

\usepackage{titlesec}
\titleformat{\section}{\Large\bfseries\color{titleblue}}{\thesection}{0.6em}{}
\titlespacing*{\section}{0pt}{1.2em}{0.5em}

\begin{document}
\sloppy

\begin{center}
{\Huge\bfseries\color{titleblue}Lektionsplan: Kraftaddition \& Kraftzerlegung}\\[0.4em]
{\large Doppellektion (2$\times$45 Min.) -- Diskussionsgrundlage}
\end{center}
\vspace{0.5em}

\noindent\textbf{Klasse:} 2e (02.09.2026) \hfill \textbf{Lehrperson:} Loris Keller

\vspace{0.3em}
\noindent\textcolor{darkgray}{Dieses Dokument ist \textbf{keine} Reinschrift des Arbeitsblatts, sondern die
Planungsgrundlage dafür: Lernziele, Aufbau der Lektion und die gewählte
Didaktik-Methode pro Phase (aus \texttt{Methodensammlung\_Physikunterricht.pdf}).
Nach Absprache wird daraus das \texttt{.tex}-Arbeitsblatt/Aufgabenblatt lokal in
VS Code erstellt. Direkte Fortsetzung von Gewichtskraft \& Federkraft
(26.08.2026) -- dieselbe Feder-Masse-Skizze dient hier als Einstieg.}

% =========================================================================
\section*{1. Abgrenzung und Lernziele}

\textbf{Im Fokus:} Kraft als vektorielle (gerichtete) Grösse einführen und
zwei grafische Grundoperationen damit üben -- die Kraftaddition (zwei und
mehrere Kräfte) und die Kraftzerlegung mit dem Kräfteparallelogramm.
Einstieg über die bereits bekannte Feder-Masse-Skizze der letzten Lektion.

\textbf{Bereits bekannt (nur Repetition, keine Neueinführung):}
Gewichtskraft $F_G=m\cdot g$ und Federkraft $F=D\cdot\Delta x$ (letzte
Doppellektion), die Feder-Masse-Skizze selbst und das Prinzip, dass sich
Feder- und Gewichtskraft im Gleichgewicht die Waage halten (dort bereits
als Kräftediagramm-Aufgabe gezeichnet, hier zum ersten Mal explizit als
Vektor-Gleichgewicht formalisiert), Kraft in Newton.

\textbf{Bewusst nicht Teil dieser Einheit:} rechnerische/komponentenweise
Kraftaddition (Sinus/Kosinus-Zerlegung) -- ausschliesslich grafisch, wie
vorgegeben. Ebenfalls nicht Teil: Normalkraft und schiefe Ebene (eigene
Doppellektion \texttt{normalkraft/}, 09.09.2026 -- das klassische
Zerlegungsbeispiel \glqq Hangabtriebskraft auf der schiefen Ebene\grqq{}
wird deshalb hier bewusst nicht verwendet, damit dort noch etwas Neues zu
entdecken bleibt), Wirklinie/Angriffspunkt-Verschiebung als eigenes
Lernziel, Drehmoment/Hebelgesetze.

\begin{itemize}[leftmargin=1.3cm,itemsep=0.3em]
  \item[LZ1] Kraft als vektorielle Grösse (Betrag \emph{und} Richtung)
  beschreiben und Kräfte an einem Körper korrekt einzeichnen (vom
  Schwerpunkt aus, mit einheitlichem Kraft-Maßstab) -- am Beispiel der
  Feder-Masse-Skizze das Kräftegleichgewicht als Vektor-Gleichgewicht
  formulieren.
  \item[LZ2] die grafische Summe (Resultierende) von zwei und von
  mehreren an einem Angriffspunkt wirkenden Kräften konstruieren
  (Kopf-an-Schwanz-Verfahren bzw. Kräfteparallelogramm bei zwei Kräften),
  mit einheitlichem Kraft-Maßstab.
  \item[LZ3] an einem selbst konstruierten Beispiel zeigen, dass der
  Betrag der resultierenden Kraft im Allgemeinen \emph{nicht} der
  arithmetischen Summe der einzelnen Kraftbeträge entspricht, sondern von
  der Winkeldifferenz der beteiligten Kräfte abhängt.
  \item[LZ4] eine gegebene Kraft mithilfe des Kräfteparallelogramms
  grafisch in zwei vorgegebene Richtungen zerlegen (Kraftzerlegung als
  Umkehrung der Kraftaddition).
  \item[LZ5] für Kraftaddition und Kraftzerlegung je mindestens ein
  Alltagsbeispiel nennen und die grafische Konstruktion darauf anwenden.
\end{itemize}

% =========================================================================
\section*{2. Aufbau der Doppellektion}

\begin{xltabular}{\textwidth}{@{}p{2.6cm}X p{1.2cm}@{}}
\toprule
\textbf{Phase} & \textbf{Inhalt \& Methode} & \textbf{Zeit} \\
\midrule
\endhead

Repetition/ Anknüpfung &
\textbf{Feder-Masse-Skizze: Kräfte einzeichnen.} Dieselbe Skizze wie in
der letzten Lernaufgabe (Feder mit hängender Masse) -- Kräfte vom
Schwerpunkt aus einzeichnen (LZ1, Repetition). \Methode{B1
Vergewisserungsphase} -- kurze Einzelarbeit, danach Plenumsabgleich. &
6' \\[0.4em]

Theorie &
\textbf{Kraft als Vektor, Kräftegleichgewicht.} Kraft hat Betrag
\emph{und} Richtung, mathematisch ein Pfeil/Vektor $\vec F$ (PhysikLibre
4.1.2). An der Feder-Masse-Skizze: Federkraft und Gewichtskraft sind
gleich lang, zeigen aber in entgegengesetzte Richtung -- Vektorsumme
gleich Null (LZ1). Lehrervortrag, geleitet, Tafelskizze. &
10' \\[0.4em]

Theorie/ Konstruktion &
\textbf{Kraftaddition zweier Kräfte, Kräfteparallelogramm, Kraft-Maßstab.}
Zwei Kräfte am selben Angriffspunkt grafisch addieren: Kräfteparallelogramm,
Resultierende als Diagonale durch den Angriffspunkt (PhysikLibre 4.3.2).
Konzept des Kraft-Maßstabs einführen (z.\,B. 1\,cm $\hat=$ 1\,N, PhysikLibre
4.3.3) -- für alle Kraftpfeile im selben Diagramm verbindlich (LZ2).
Lehrervortrag, gemeinsame Tafelkonstruktion. &
10' \\[0.4em]

Übung &
\textbf{Kraftaddition mehrerer Kräfte, Alltagsbeispiele.} Kopf-an-Schwanz-
Verfahren für mehr als zwei Kräfte am selben Angriffspunkt (PhysikLibre
4.3.1) -- Reihenfolge der Pfeile ist beliebig, Ergebnis bleibt gleich.
Gruppen konstruieren je ein eigenes Alltagsbeispiel (Tauziehen, Boot mit
Motor und Seitenwind, Hund an zwei Leinen) auf dem Whiteboard und
präsentieren kurz (LZ2, LZ5). \Methode{M3 Whiteboarding} -- Dreiergruppen,
danach Boards hochhalten. &
15' \\[0.4em]

Kernaufgabe &
\textbf{Resultierende $\neq$ Summe der Beträge.} Zwei gleich grosse Kräfte
(z.\,B. je 4\,N) bei verschiedenen Zwischenwinkeln (0°, 90°, 180°, ggf.
60°) grafisch addieren, Resultierende jeweils mit dem Maßstab ausmessen.
Entdeckung: bei 0° ist die Resultierende maximal (Summe der Beträge), bei
180° minimal, dazwischen liegt sie immer \emph{unterhalb} der
arithmetischen Summe (LZ3). \Methode{M9 Produktives Üben} -- Winkel
systematisch variieren, Muster selbst entdecken; Einzel-/Partnerarbeit. &
15' \\[0.4em]

Theorie &
\textbf{Kraftzerlegung mit dem Kräfteparallelogramm.} Umkehrung von
Kraftaddition: eine gegebene Kraft in zwei vorgegebene Richtungen
zerlegen (PhysikLibre 4.3.4). Alltagsbeispiel: eine Verkehrsampel, die an
zwei schräg gespannten Seilen hängt -- die Gewichtskraft wird in zwei
Teilkräfte entlang der Seilrichtungen zerlegt (Bild 4.17) (LZ4).
Lehrervortrag, geleitet, Tafelkonstruktion. &
10' \\[0.4em]

Übung &
\textbf{Kraftzerlegung selbst konstruieren.} Partnerarbeit: eigene
Zerlegungsaufgabe konstruieren (vorgegebene Kraft, vorgegebene Richtungen),
plus ein selbst gewähltes Alltagsbeispiel (Zerlegung \emph{oder} weitere
Addition) grafisch umsetzen (LZ4, LZ5). Alle Konstruktionen auf Feinraster
mit angegebenem Maßstab, Resultat mit Lineal nachmessbar. Partnerarbeit,
Lehrperson zirkuliert. &
15' \\[0.4em]

Abschluss &
\textbf{Zusammenfassung und Ausblick aufs Aufgabenblatt.} Kurze
Zusammenfassung der beiden Grundoperationen (Addition/Zerlegung) und
Hinweis auf das separate Aufgabenblatt für weitere Übung. Plenum. &
5' \\

\bottomrule
\end{xltabular}

\textcolor{darkgray}{\small Summe: 86 Min. -- 4 Min. Puffer innerhalb der
90 Min. der Doppellektion, bewusst als Reserve für die beiden
Konstruktionsphasen (Whiteboarding, Kernaufgabe), deren Zeitbedarf je nach
Zeichentempo der Klasse schwankt.}

% =========================================================================
\section*{3. Begründung der Methodenwahl}

Die Kernaufgabe (Resultierende $\neq$ Summe der Beträge) ist bewusst als
M9 Produktives Üben statt als einzelne Rechenaufgabe angelegt: Durch
systematisches Variieren des Zwischenwinkels bei sonst gleichbleibenden
Kraftbeträgen entdecken die SuS das Muster (Maximum bei 0°, Minimum bei
180°) selbst, statt es nur vorgetragen zu bekommen -- passend zur
expliziten Vorgabe, dass mindestens ein Beispiel zeigen soll, dass Kräfte
sich nicht wie gewöhnliche Zahlen addieren.

Die Alltagsbeispiel-Phase (Kraftaddition mehrerer Kräfte) nutzt M3
Whiteboarding statt Lehrervortrag, weil die Konstruktion selbst
(Kopf-an-Schwanz-Verfahren) beim eigenen Zeichnen viel eher verinnerlicht
wird als beim Zusehen, und weil Alltagsbeispiele in der Gruppe
diskutiert und verglichen werden sollen (mehrere plausible Beispiele
kursieren erfahrungsgemäss, z.\,B. Tauziehen vs. Boot mit Seitenwind).

Die abschliessende Zerlegungs-Übung ist bewusst Partnerarbeit statt
Whiteboarding: Nach zwei bereits kollaborativen Phasen (Addition-Übung,
Kernaufgabe) sorgt der Wechsel zur Partnerarbeit für Abwechslung und gibt
jeder/jedem genügend eigene Zeichenzeit für eine individuell wählbare
Aufgabe.

Der Einstieg über die Feder-Masse-Skizze knüpft bewusst an eine bereits
gezeichnete, vertraute Situation an (Sara Romers Vorgabe: Anknüpfung an
die letzte Lektion), statt mit einem neuen, unbekannten Beispiel zu
beginnen -- das Kräftegleichgewicht wird dadurch nicht neu eingeführt,
sondern nur neu (als Vektorgleichung) interpretiert.

% =========================================================================
\section*{4. Grafische Aufgaben: Feinraster und Maßstab (verbindlich)}

Wie vorgegeben erhält \emph{jede} grafische Konstruktionsaufgabe --
sowohl im Arbeitsblatt als auch im Aufgabenblatt -- ein Feld mit feinem
Punkt- oder Liniengitter sowie eine explizit angegebene Maßstabsangabe
(z.\,B. \glqq 1\,cm $\hat=$ 1\,N\grqq{}). Die gegebenen Kraftpfeile sind
exakt in der angegebenen Länge einzuzeichnen, sodass auch das Ergebnis
(Resultierende bzw. Teilkräfte) mit demselben Maßstab nachgemessen werden
kann -- die SuS sollen ihre Konstruktion selbst kontrollieren können,
ohne die Lösung nachschlagen zu müssen. Diese Vorgabe gilt für alle
Konstruktionsaufgaben der Kernaufgabe, der Whiteboarding-Übung und der
abschliessenden Partnerarbeit gleichermassen.

% =========================================================================
\section*{5. Anschlussdokument: Aufgabenblatt}

Wie von Sara Romer gewünscht, entsteht zusätzlich ein separates
Aufgabenblatt mit reinen Übungsaufgaben zu Kraftaddition und
-zerlegung -- alle grafisch, mit Feinraster und Maßstab wie oben --
inklusive einer Aufgabe, in der die SuS ein eigenes Alltagsbeispiel frei
wählen (Addition \emph{oder} Zerlegung) und selbst konstruieren.

% =========================================================================
\section*{6. Optionale Zusatzmaterialien (Selbstkontrolle, nicht Ersatz)}

Als optionale Ergänzung zur eigenen grafischen Konstruktion, nicht als
Ersatz dafür: die PhET-Simulationen \glqq Vector Addition\grqq{} (grafisch,
Kopf-an-Schwanz und Parallelogramm) sowie \glqq Vector Addition:
Equations\grqq{} (Komponentenschreibweise -- bewusst nicht Teil dieser
Lektion, siehe Abschnitt 1).

\end{document}
